\documentclass[11pt]{article}

    \usepackage[breakable]{tcolorbox}
    \usepackage{parskip} % Stop auto-indenting (to mimic markdown behaviour)
    

    % Basic figure setup, for now with no caption control since it's done
    % automatically by Pandoc (which extracts ![](path) syntax from Markdown).
    \usepackage{graphicx}
    % Maintain compatibility with old templates. Remove in nbconvert 6.0
    \let\Oldincludegraphics\includegraphics
    % Ensure that by default, figures have no caption (until we provide a
    % proper Figure object with a Caption API and a way to capture that
    % in the conversion process - todo).
    \usepackage{caption}
    \DeclareCaptionFormat{nocaption}{}
    \captionsetup{format=nocaption,aboveskip=0pt,belowskip=0pt}

    \usepackage{float}
    \floatplacement{figure}{H} % forces figures to be placed at the correct location
    \usepackage{xcolor} % Allow colors to be defined
    \usepackage{enumerate} % Needed for markdown enumerations to work
    \usepackage{geometry} % Used to adjust the document margins
    \usepackage{amsmath} % Equations
    \usepackage{amssymb} % Equations
    \usepackage{textcomp} % defines textquotesingle
    % Hack from http://tex.stackexchange.com/a/47451/13684:
    \AtBeginDocument{%
        \def\PYZsq{\textquotesingle}% Upright quotes in Pygmentized code
    }
    \usepackage{upquote} % Upright quotes for verbatim code
    \usepackage{eurosym} % defines \euro

    \usepackage{iftex}
    \ifPDFTeX
        \usepackage[T1]{fontenc}
        \IfFileExists{alphabeta.sty}{
              \usepackage{alphabeta}
          }{
              \usepackage[mathletters]{ucs}
              \usepackage[utf8x]{inputenc}
          }
    \else
        \usepackage{fontspec}
        \usepackage{unicode-math}
    \fi

    \usepackage{fancyvrb} % verbatim replacement that allows latex
    \usepackage{grffile} % extends the file name processing of package graphics
                         % to support a larger range
    \makeatletter % fix for old versions of grffile with XeLaTeX
    \@ifpackagelater{grffile}{2019/11/01}
    {
      % Do nothing on new versions
    }
    {
      \def\Gread@@xetex#1{%
        \IfFileExists{"\Gin@base".bb}%
        {\Gread@eps{\Gin@base.bb}}%
        {\Gread@@xetex@aux#1}%
      }
    }
    \makeatother
    \usepackage[Export]{adjustbox} % Used to constrain images to a maximum size
    \adjustboxset{max size={0.9\linewidth}{0.9\paperheight}}

    % The hyperref package gives us a pdf with properly built
    % internal navigation ('pdf bookmarks' for the table of contents,
    % internal cross-reference links, web links for URLs, etc.)
    \usepackage{hyperref}
    % The default LaTeX title has an obnoxious amount of whitespace. By default,
    % titling removes some of it. It also provides customization options.
    \usepackage{titling}
    \usepackage{longtable} % longtable support required by pandoc >1.10
    \usepackage{booktabs}  % table support for pandoc > 1.12.2
    \usepackage{array}     % table support for pandoc >= 2.11.3
    \usepackage{calc}      % table minipage width calculation for pandoc >= 2.11.1
    \usepackage[inline]{enumitem} % IRkernel/repr support (it uses the enumerate* environment)
    \usepackage[normalem]{ulem} % ulem is needed to support strikethroughs (\sout)
                                % normalem makes italics be italics, not underlines
    \usepackage{soul}      % strikethrough (\st) support for pandoc >= 3.0.0
    \usepackage{mathrsfs}
    

    
    % Colors for the hyperref package
    \definecolor{urlcolor}{rgb}{0,.145,.698}
    \definecolor{linkcolor}{rgb}{.71,0.21,0.01}
    \definecolor{citecolor}{rgb}{.12,.54,.11}

    % ANSI colors
    \definecolor{ansi-black}{HTML}{3E424D}
    \definecolor{ansi-black-intense}{HTML}{282C36}
    \definecolor{ansi-red}{HTML}{E75C58}
    \definecolor{ansi-red-intense}{HTML}{B22B31}
    \definecolor{ansi-green}{HTML}{00A250}
    \definecolor{ansi-green-intense}{HTML}{007427}
    \definecolor{ansi-yellow}{HTML}{DDB62B}
    \definecolor{ansi-yellow-intense}{HTML}{B27D12}
    \definecolor{ansi-blue}{HTML}{208FFB}
    \definecolor{ansi-blue-intense}{HTML}{0065CA}
    \definecolor{ansi-magenta}{HTML}{D160C4}
    \definecolor{ansi-magenta-intense}{HTML}{A03196}
    \definecolor{ansi-cyan}{HTML}{60C6C8}
    \definecolor{ansi-cyan-intense}{HTML}{258F8F}
    \definecolor{ansi-white}{HTML}{C5C1B4}
    \definecolor{ansi-white-intense}{HTML}{A1A6B2}
    \definecolor{ansi-default-inverse-fg}{HTML}{FFFFFF}
    \definecolor{ansi-default-inverse-bg}{HTML}{000000}

    % common color for the border for error outputs.
    \definecolor{outerrorbackground}{HTML}{FFDFDF}

    % commands and environments needed by pandoc snippets
    % extracted from the output of `pandoc -s`
    \providecommand{\tightlist}{%
      \setlength{\itemsep}{0pt}\setlength{\parskip}{0pt}}
    \DefineVerbatimEnvironment{Highlighting}{Verbatim}{commandchars=\\\{\}}
    % Add ',fontsize=\small' for more characters per line
    \newenvironment{Shaded}{}{}
    \newcommand{\KeywordTok}[1]{\textcolor[rgb]{0.00,0.44,0.13}{\textbf{{#1}}}}
    \newcommand{\DataTypeTok}[1]{\textcolor[rgb]{0.56,0.13,0.00}{{#1}}}
    \newcommand{\DecValTok}[1]{\textcolor[rgb]{0.25,0.63,0.44}{{#1}}}
    \newcommand{\BaseNTok}[1]{\textcolor[rgb]{0.25,0.63,0.44}{{#1}}}
    \newcommand{\FloatTok}[1]{\textcolor[rgb]{0.25,0.63,0.44}{{#1}}}
    \newcommand{\CharTok}[1]{\textcolor[rgb]{0.25,0.44,0.63}{{#1}}}
    \newcommand{\StringTok}[1]{\textcolor[rgb]{0.25,0.44,0.63}{{#1}}}
    \newcommand{\CommentTok}[1]{\textcolor[rgb]{0.38,0.63,0.69}{\textit{{#1}}}}
    \newcommand{\OtherTok}[1]{\textcolor[rgb]{0.00,0.44,0.13}{{#1}}}
    \newcommand{\AlertTok}[1]{\textcolor[rgb]{1.00,0.00,0.00}{\textbf{{#1}}}}
    \newcommand{\FunctionTok}[1]{\textcolor[rgb]{0.02,0.16,0.49}{{#1}}}
    \newcommand{\RegionMarkerTok}[1]{{#1}}
    \newcommand{\ErrorTok}[1]{\textcolor[rgb]{1.00,0.00,0.00}{\textbf{{#1}}}}
    \newcommand{\NormalTok}[1]{{#1}}

    % Additional commands for more recent versions of Pandoc
    \newcommand{\ConstantTok}[1]{\textcolor[rgb]{0.53,0.00,0.00}{{#1}}}
    \newcommand{\SpecialCharTok}[1]{\textcolor[rgb]{0.25,0.44,0.63}{{#1}}}
    \newcommand{\VerbatimStringTok}[1]{\textcolor[rgb]{0.25,0.44,0.63}{{#1}}}
    \newcommand{\SpecialStringTok}[1]{\textcolor[rgb]{0.73,0.40,0.53}{{#1}}}
    \newcommand{\ImportTok}[1]{{#1}}
    \newcommand{\DocumentationTok}[1]{\textcolor[rgb]{0.73,0.13,0.13}{\textit{{#1}}}}
    \newcommand{\AnnotationTok}[1]{\textcolor[rgb]{0.38,0.63,0.69}{\textbf{\textit{{#1}}}}}
    \newcommand{\CommentVarTok}[1]{\textcolor[rgb]{0.38,0.63,0.69}{\textbf{\textit{{#1}}}}}
    \newcommand{\VariableTok}[1]{\textcolor[rgb]{0.10,0.09,0.49}{{#1}}}
    \newcommand{\ControlFlowTok}[1]{\textcolor[rgb]{0.00,0.44,0.13}{\textbf{{#1}}}}
    \newcommand{\OperatorTok}[1]{\textcolor[rgb]{0.40,0.40,0.40}{{#1}}}
    \newcommand{\BuiltInTok}[1]{{#1}}
    \newcommand{\ExtensionTok}[1]{{#1}}
    \newcommand{\PreprocessorTok}[1]{\textcolor[rgb]{0.74,0.48,0.00}{{#1}}}
    \newcommand{\AttributeTok}[1]{\textcolor[rgb]{0.49,0.56,0.16}{{#1}}}
    \newcommand{\InformationTok}[1]{\textcolor[rgb]{0.38,0.63,0.69}{\textbf{\textit{{#1}}}}}
    \newcommand{\WarningTok}[1]{\textcolor[rgb]{0.38,0.63,0.69}{\textbf{\textit{{#1}}}}}


    % Define a nice break command that doesn't care if a line doesn't already
    % exist.
    \def\br{\hspace*{\fill} \\* }
    % Math Jax compatibility definitions
    \def\gt{>}
    \def\lt{<}
    \let\Oldtex\TeX
    \let\Oldlatex\LaTeX
    \renewcommand{\TeX}{\textrm{\Oldtex}}
    \renewcommand{\LaTeX}{\textrm{\Oldlatex}}
    % Document parameters
    % Document title
    \title{exam\_storyline}
    
    
    
    
    
    
    
% Pygments definitions
\makeatletter
\def\PY@reset{\let\PY@it=\relax \let\PY@bf=\relax%
    \let\PY@ul=\relax \let\PY@tc=\relax%
    \let\PY@bc=\relax \let\PY@ff=\relax}
\def\PY@tok#1{\csname PY@tok@#1\endcsname}
\def\PY@toks#1+{\ifx\relax#1\empty\else%
    \PY@tok{#1}\expandafter\PY@toks\fi}
\def\PY@do#1{\PY@bc{\PY@tc{\PY@ul{%
    \PY@it{\PY@bf{\PY@ff{#1}}}}}}}
\def\PY#1#2{\PY@reset\PY@toks#1+\relax+\PY@do{#2}}

\@namedef{PY@tok@w}{\def\PY@tc##1{\textcolor[rgb]{0.73,0.73,0.73}{##1}}}
\@namedef{PY@tok@c}{\let\PY@it=\textit\def\PY@tc##1{\textcolor[rgb]{0.24,0.48,0.48}{##1}}}
\@namedef{PY@tok@cp}{\def\PY@tc##1{\textcolor[rgb]{0.61,0.40,0.00}{##1}}}
\@namedef{PY@tok@k}{\let\PY@bf=\textbf\def\PY@tc##1{\textcolor[rgb]{0.00,0.50,0.00}{##1}}}
\@namedef{PY@tok@kp}{\def\PY@tc##1{\textcolor[rgb]{0.00,0.50,0.00}{##1}}}
\@namedef{PY@tok@kt}{\def\PY@tc##1{\textcolor[rgb]{0.69,0.00,0.25}{##1}}}
\@namedef{PY@tok@o}{\def\PY@tc##1{\textcolor[rgb]{0.40,0.40,0.40}{##1}}}
\@namedef{PY@tok@ow}{\let\PY@bf=\textbf\def\PY@tc##1{\textcolor[rgb]{0.67,0.13,1.00}{##1}}}
\@namedef{PY@tok@nb}{\def\PY@tc##1{\textcolor[rgb]{0.00,0.50,0.00}{##1}}}
\@namedef{PY@tok@nf}{\def\PY@tc##1{\textcolor[rgb]{0.00,0.00,1.00}{##1}}}
\@namedef{PY@tok@nc}{\let\PY@bf=\textbf\def\PY@tc##1{\textcolor[rgb]{0.00,0.00,1.00}{##1}}}
\@namedef{PY@tok@nn}{\let\PY@bf=\textbf\def\PY@tc##1{\textcolor[rgb]{0.00,0.00,1.00}{##1}}}
\@namedef{PY@tok@ne}{\let\PY@bf=\textbf\def\PY@tc##1{\textcolor[rgb]{0.80,0.25,0.22}{##1}}}
\@namedef{PY@tok@nv}{\def\PY@tc##1{\textcolor[rgb]{0.10,0.09,0.49}{##1}}}
\@namedef{PY@tok@no}{\def\PY@tc##1{\textcolor[rgb]{0.53,0.00,0.00}{##1}}}
\@namedef{PY@tok@nl}{\def\PY@tc##1{\textcolor[rgb]{0.46,0.46,0.00}{##1}}}
\@namedef{PY@tok@ni}{\let\PY@bf=\textbf\def\PY@tc##1{\textcolor[rgb]{0.44,0.44,0.44}{##1}}}
\@namedef{PY@tok@na}{\def\PY@tc##1{\textcolor[rgb]{0.41,0.47,0.13}{##1}}}
\@namedef{PY@tok@nt}{\let\PY@bf=\textbf\def\PY@tc##1{\textcolor[rgb]{0.00,0.50,0.00}{##1}}}
\@namedef{PY@tok@nd}{\def\PY@tc##1{\textcolor[rgb]{0.67,0.13,1.00}{##1}}}
\@namedef{PY@tok@s}{\def\PY@tc##1{\textcolor[rgb]{0.73,0.13,0.13}{##1}}}
\@namedef{PY@tok@sd}{\let\PY@it=\textit\def\PY@tc##1{\textcolor[rgb]{0.73,0.13,0.13}{##1}}}
\@namedef{PY@tok@si}{\let\PY@bf=\textbf\def\PY@tc##1{\textcolor[rgb]{0.64,0.35,0.47}{##1}}}
\@namedef{PY@tok@se}{\let\PY@bf=\textbf\def\PY@tc##1{\textcolor[rgb]{0.67,0.36,0.12}{##1}}}
\@namedef{PY@tok@sr}{\def\PY@tc##1{\textcolor[rgb]{0.64,0.35,0.47}{##1}}}
\@namedef{PY@tok@ss}{\def\PY@tc##1{\textcolor[rgb]{0.10,0.09,0.49}{##1}}}
\@namedef{PY@tok@sx}{\def\PY@tc##1{\textcolor[rgb]{0.00,0.50,0.00}{##1}}}
\@namedef{PY@tok@m}{\def\PY@tc##1{\textcolor[rgb]{0.40,0.40,0.40}{##1}}}
\@namedef{PY@tok@gh}{\let\PY@bf=\textbf\def\PY@tc##1{\textcolor[rgb]{0.00,0.00,0.50}{##1}}}
\@namedef{PY@tok@gu}{\let\PY@bf=\textbf\def\PY@tc##1{\textcolor[rgb]{0.50,0.00,0.50}{##1}}}
\@namedef{PY@tok@gd}{\def\PY@tc##1{\textcolor[rgb]{0.63,0.00,0.00}{##1}}}
\@namedef{PY@tok@gi}{\def\PY@tc##1{\textcolor[rgb]{0.00,0.52,0.00}{##1}}}
\@namedef{PY@tok@gr}{\def\PY@tc##1{\textcolor[rgb]{0.89,0.00,0.00}{##1}}}
\@namedef{PY@tok@ge}{\let\PY@it=\textit}
\@namedef{PY@tok@gs}{\let\PY@bf=\textbf}
\@namedef{PY@tok@ges}{\let\PY@bf=\textbf\let\PY@it=\textit}
\@namedef{PY@tok@gp}{\let\PY@bf=\textbf\def\PY@tc##1{\textcolor[rgb]{0.00,0.00,0.50}{##1}}}
\@namedef{PY@tok@go}{\def\PY@tc##1{\textcolor[rgb]{0.44,0.44,0.44}{##1}}}
\@namedef{PY@tok@gt}{\def\PY@tc##1{\textcolor[rgb]{0.00,0.27,0.87}{##1}}}
\@namedef{PY@tok@err}{\def\PY@bc##1{{\setlength{\fboxsep}{\string -\fboxrule}\fcolorbox[rgb]{1.00,0.00,0.00}{1,1,1}{\strut ##1}}}}
\@namedef{PY@tok@kc}{\let\PY@bf=\textbf\def\PY@tc##1{\textcolor[rgb]{0.00,0.50,0.00}{##1}}}
\@namedef{PY@tok@kd}{\let\PY@bf=\textbf\def\PY@tc##1{\textcolor[rgb]{0.00,0.50,0.00}{##1}}}
\@namedef{PY@tok@kn}{\let\PY@bf=\textbf\def\PY@tc##1{\textcolor[rgb]{0.00,0.50,0.00}{##1}}}
\@namedef{PY@tok@kr}{\let\PY@bf=\textbf\def\PY@tc##1{\textcolor[rgb]{0.00,0.50,0.00}{##1}}}
\@namedef{PY@tok@bp}{\def\PY@tc##1{\textcolor[rgb]{0.00,0.50,0.00}{##1}}}
\@namedef{PY@tok@fm}{\def\PY@tc##1{\textcolor[rgb]{0.00,0.00,1.00}{##1}}}
\@namedef{PY@tok@vc}{\def\PY@tc##1{\textcolor[rgb]{0.10,0.09,0.49}{##1}}}
\@namedef{PY@tok@vg}{\def\PY@tc##1{\textcolor[rgb]{0.10,0.09,0.49}{##1}}}
\@namedef{PY@tok@vi}{\def\PY@tc##1{\textcolor[rgb]{0.10,0.09,0.49}{##1}}}
\@namedef{PY@tok@vm}{\def\PY@tc##1{\textcolor[rgb]{0.10,0.09,0.49}{##1}}}
\@namedef{PY@tok@sa}{\def\PY@tc##1{\textcolor[rgb]{0.73,0.13,0.13}{##1}}}
\@namedef{PY@tok@sb}{\def\PY@tc##1{\textcolor[rgb]{0.73,0.13,0.13}{##1}}}
\@namedef{PY@tok@sc}{\def\PY@tc##1{\textcolor[rgb]{0.73,0.13,0.13}{##1}}}
\@namedef{PY@tok@dl}{\def\PY@tc##1{\textcolor[rgb]{0.73,0.13,0.13}{##1}}}
\@namedef{PY@tok@s2}{\def\PY@tc##1{\textcolor[rgb]{0.73,0.13,0.13}{##1}}}
\@namedef{PY@tok@sh}{\def\PY@tc##1{\textcolor[rgb]{0.73,0.13,0.13}{##1}}}
\@namedef{PY@tok@s1}{\def\PY@tc##1{\textcolor[rgb]{0.73,0.13,0.13}{##1}}}
\@namedef{PY@tok@mb}{\def\PY@tc##1{\textcolor[rgb]{0.40,0.40,0.40}{##1}}}
\@namedef{PY@tok@mf}{\def\PY@tc##1{\textcolor[rgb]{0.40,0.40,0.40}{##1}}}
\@namedef{PY@tok@mh}{\def\PY@tc##1{\textcolor[rgb]{0.40,0.40,0.40}{##1}}}
\@namedef{PY@tok@mi}{\def\PY@tc##1{\textcolor[rgb]{0.40,0.40,0.40}{##1}}}
\@namedef{PY@tok@il}{\def\PY@tc##1{\textcolor[rgb]{0.40,0.40,0.40}{##1}}}
\@namedef{PY@tok@mo}{\def\PY@tc##1{\textcolor[rgb]{0.40,0.40,0.40}{##1}}}
\@namedef{PY@tok@ch}{\let\PY@it=\textit\def\PY@tc##1{\textcolor[rgb]{0.24,0.48,0.48}{##1}}}
\@namedef{PY@tok@cm}{\let\PY@it=\textit\def\PY@tc##1{\textcolor[rgb]{0.24,0.48,0.48}{##1}}}
\@namedef{PY@tok@cpf}{\let\PY@it=\textit\def\PY@tc##1{\textcolor[rgb]{0.24,0.48,0.48}{##1}}}
\@namedef{PY@tok@c1}{\let\PY@it=\textit\def\PY@tc##1{\textcolor[rgb]{0.24,0.48,0.48}{##1}}}
\@namedef{PY@tok@cs}{\let\PY@it=\textit\def\PY@tc##1{\textcolor[rgb]{0.24,0.48,0.48}{##1}}}

\def\PYZbs{\char`\\}
\def\PYZus{\char`\_}
\def\PYZob{\char`\{}
\def\PYZcb{\char`\}}
\def\PYZca{\char`\^}
\def\PYZam{\char`\&}
\def\PYZlt{\char`\<}
\def\PYZgt{\char`\>}
\def\PYZsh{\char`\#}
\def\PYZpc{\char`\%}
\def\PYZdl{\char`\$}
\def\PYZhy{\char`\-}
\def\PYZsq{\char`\'}
\def\PYZdq{\char`\"}
\def\PYZti{\char`\~}
% for compatibility with earlier versions
\def\PYZat{@}
\def\PYZlb{[}
\def\PYZrb{]}
\makeatother


    % For linebreaks inside Verbatim environment from package fancyvrb.
    \makeatletter
        \newbox\Wrappedcontinuationbox
        \newbox\Wrappedvisiblespacebox
        \newcommand*\Wrappedvisiblespace {\textcolor{red}{\textvisiblespace}}
        \newcommand*\Wrappedcontinuationsymbol {\textcolor{red}{\llap{\tiny$\m@th\hookrightarrow$}}}
        \newcommand*\Wrappedcontinuationindent {3ex }
        \newcommand*\Wrappedafterbreak {\kern\Wrappedcontinuationindent\copy\Wrappedcontinuationbox}
        % Take advantage of the already applied Pygments mark-up to insert
        % potential linebreaks for TeX processing.
        %        {, <, #, %, $, ' and ": go to next line.
        %        _, }, ^, &, >, - and ~: stay at end of broken line.
        % Use of \textquotesingle for straight quote.
        \newcommand*\Wrappedbreaksatspecials {%
            \def\PYGZus{\discretionary{\char`\_}{\Wrappedafterbreak}{\char`\_}}%
            \def\PYGZob{\discretionary{}{\Wrappedafterbreak\char`\{}{\char`\{}}%
            \def\PYGZcb{\discretionary{\char`\}}{\Wrappedafterbreak}{\char`\}}}%
            \def\PYGZca{\discretionary{\char`\^}{\Wrappedafterbreak}{\char`\^}}%
            \def\PYGZam{\discretionary{\char`\&}{\Wrappedafterbreak}{\char`\&}}%
            \def\PYGZlt{\discretionary{}{\Wrappedafterbreak\char`\<}{\char`\<}}%
            \def\PYGZgt{\discretionary{\char`\>}{\Wrappedafterbreak}{\char`\>}}%
            \def\PYGZsh{\discretionary{}{\Wrappedafterbreak\char`\#}{\char`\#}}%
            \def\PYGZpc{\discretionary{}{\Wrappedafterbreak\char`\%}{\char`\%}}%
            \def\PYGZdl{\discretionary{}{\Wrappedafterbreak\char`\$}{\char`\$}}%
            \def\PYGZhy{\discretionary{\char`\-}{\Wrappedafterbreak}{\char`\-}}%
            \def\PYGZsq{\discretionary{}{\Wrappedafterbreak\textquotesingle}{\textquotesingle}}%
            \def\PYGZdq{\discretionary{}{\Wrappedafterbreak\char`\"}{\char`\"}}%
            \def\PYGZti{\discretionary{\char`\~}{\Wrappedafterbreak}{\char`\~}}%
        }
        % Some characters . , ; ? ! / are not pygmentized.
        % This macro makes them "active" and they will insert potential linebreaks
        \newcommand*\Wrappedbreaksatpunct {%
            \lccode`\~`\.\lowercase{\def~}{\discretionary{\hbox{\char`\.}}{\Wrappedafterbreak}{\hbox{\char`\.}}}%
            \lccode`\~`\,\lowercase{\def~}{\discretionary{\hbox{\char`\,}}{\Wrappedafterbreak}{\hbox{\char`\,}}}%
            \lccode`\~`\;\lowercase{\def~}{\discretionary{\hbox{\char`\;}}{\Wrappedafterbreak}{\hbox{\char`\;}}}%
            \lccode`\~`\:\lowercase{\def~}{\discretionary{\hbox{\char`\:}}{\Wrappedafterbreak}{\hbox{\char`\:}}}%
            \lccode`\~`\?\lowercase{\def~}{\discretionary{\hbox{\char`\?}}{\Wrappedafterbreak}{\hbox{\char`\?}}}%
            \lccode`\~`\!\lowercase{\def~}{\discretionary{\hbox{\char`\!}}{\Wrappedafterbreak}{\hbox{\char`\!}}}%
            \lccode`\~`\/\lowercase{\def~}{\discretionary{\hbox{\char`\/}}{\Wrappedafterbreak}{\hbox{\char`\/}}}%
            \catcode`\.\active
            \catcode`\,\active
            \catcode`\;\active
            \catcode`\:\active
            \catcode`\?\active
            \catcode`\!\active
            \catcode`\/\active
            \lccode`\~`\~
        }
    \makeatother

    \let\OriginalVerbatim=\Verbatim
    \makeatletter
    \renewcommand{\Verbatim}[1][1]{%
        %\parskip\z@skip
        \sbox\Wrappedcontinuationbox {\Wrappedcontinuationsymbol}%
        \sbox\Wrappedvisiblespacebox {\FV@SetupFont\Wrappedvisiblespace}%
        \def\FancyVerbFormatLine ##1{\hsize\linewidth
            \vtop{\raggedright\hyphenpenalty\z@\exhyphenpenalty\z@
                \doublehyphendemerits\z@\finalhyphendemerits\z@
                \strut ##1\strut}%
        }%
        % If the linebreak is at a space, the latter will be displayed as visible
        % space at end of first line, and a continuation symbol starts next line.
        % Stretch/shrink are however usually zero for typewriter font.
        \def\FV@Space {%
            \nobreak\hskip\z@ plus\fontdimen3\font minus\fontdimen4\font
            \discretionary{\copy\Wrappedvisiblespacebox}{\Wrappedafterbreak}
            {\kern\fontdimen2\font}%
        }%

        % Allow breaks at special characters using \PYG... macros.
        \Wrappedbreaksatspecials
        % Breaks at punctuation characters . , ; ? ! and / need catcode=\active
        \OriginalVerbatim[#1,codes*=\Wrappedbreaksatpunct]%
    }
    \makeatother

    % Exact colors from NB
    \definecolor{incolor}{HTML}{303F9F}
    \definecolor{outcolor}{HTML}{D84315}
    \definecolor{cellborder}{HTML}{CFCFCF}
    \definecolor{cellbackground}{HTML}{F7F7F7}

    % prompt
    \makeatletter
    \newcommand{\boxspacing}{\kern\kvtcb@left@rule\kern\kvtcb@boxsep}
    \makeatother
    \newcommand{\prompt}[4]{
        {\ttfamily\llap{{\color{#2}[#3]:\hspace{3pt}#4}}\vspace{-\baselineskip}}
    }
    

    
    % Prevent overflowing lines due to hard-to-break entities
    \sloppy
    % Setup hyperref package
    \hypersetup{
      breaklinks=true,  % so long urls are correctly broken across lines
      colorlinks=true,
      urlcolor=urlcolor,
      linkcolor=linkcolor,
      citecolor=citecolor,
      }
    % Slightly bigger margins than the latex defaults
    
    \geometry{verbose,tmargin=1in,bmargin=1in,lmargin=1in,rmargin=1in}
    
    

\begin{document}
    
    \maketitle
    
    

    
    \section{Resilient EV Charging Under
Disruption}\label{resilient-ev-charging-under-disruption}

\subsection{Technical report notebook for DTU 42578 Advanced Business
Analytics}\label{technical-report-notebook-for-dtu-42578-advanced-business-analytics}

This notebook is the final technical report for the project on
resilience in EV charging operations. It is written as a structured
report rather than an exploration log. The central question is whether
mobile charging stations can be used intelligently to reduce queueing
pressure when demand is uneven and charging infrastructure is disrupted.

The intended research question is:

\textbf{Can a reinforcement learning policy improve resilience in an EV
charging corridor by dynamically reallocating mobile charging stations
to reduce queue length and waiting time under disruptions?}

A crucial caveat shapes the report: the current checkout contains the
active simulator, the active configs, copied RL training history, and
baseline rollout artifacts, but it does \textbf{not} contain a
current-compatible RL checkpoint or RL rollout artifact for the active
27-feature, 5-action A-B-C benchmark. The report therefore separates
three things carefully:

\begin{itemize}
\tightlist
\item
  the \textbf{RL formulation}, which can be described precisely,
\item
  the \textbf{RL training diagnostics}, which are available from copied
  training history,
\item
  the \textbf{reproducible operational rollout evidence}, which is
  currently available for the fixed baseline and heuristic adaptive
  baselines.
\end{itemize}

The notebook follows this exact structure:

\begin{enumerate}
\def\labelenumi{\arabic{enumi}.}
\tightlist
\item
  Motivation / problem definition\\
\item
  Simulation environment\\
\item
  Simulation design\\
\item
  RL agent and methods\\
\item
  Reward calibration\\
\item
  Training setup\\
\item
  Training the model\\
\item
  Simulation comparison rollout\\
\item
  Summary metrics table\\
\item
  Spatial awareness\\
\item
  Limitations\\
\item
  Conclusion\\
\item
  Appendix
\end{enumerate}

    \subsection{1. Motivation / problem
definition}\label{motivation-problem-definition}

Fast-charging infrastructure for electric vehicles is expensive,
geographically fixed, and often unevenly used over the day. In practice,
this creates four operational problems at the same time:

\begin{itemize}
\tightlist
\item
  charging demand is not constant, because traffic has strong
  time-of-day peaks,
\item
  congestion is local, because one station can become stressed while
  another still has spare capacity,
\item
  disruptions can suddenly reduce usable capacity or sharply increase
  demand,
\item
  building permanent spare infrastructure everywhere is costly and can
  leave assets under-used most of the time.
\end{itemize}

The project studies a corridor setting in which these problems are
deliberately combined. The practical goal is simple:

\textbf{use mobile charging stations intelligently to reduce queue
length and waiting time when pressure shifts across the corridor or when
a disruption occurs.}

That is a resilience question. The interest is not only average
performance on a normal day, but whether the charging system can keep
serving vehicles when conditions deteriorate.

    \subsection{2. Simulation environment}\label{simulation-environment}

The active benchmark is a synthetic A-B-C corridor with three cities
arranged on a line and two fixed fast-charging stations:

\begin{itemize}
\tightlist
\item
  one fixed station between city A and city B,
\item
  one fixed station between city B and city C,
\item
  a middle depot at city B from which mobile charging stations (MCS) can
  be dispatched.
\end{itemize}

The simulator runs in 5-minute decision steps. Vehicles enter the
corridor, some decide to charge, they queue if no plug is free, and they
start service once capacity becomes available. Mobile charging stations
add temporary extra chargers, but they are not free to teleport: they
must travel from the depot, can be moved between stations, and need time
to recharge before they can be reused.

A useful operational calibration is how much one MCS can actually do. In
the active environment, one MCS has two chargers and the average service
time is 30 minutes. That means one MCS can complete roughly \textbf{4
charging sessions per hour}, or about \textbf{0.33 vehicles per 5-minute
step on average}. This makes MCS deployment meaningful but not magically
dominant.

    \begin{longtable}[]{@{}
  >{\raggedright\arraybackslash}p{(\columnwidth - 2\tabcolsep) * \real{0.5000}}
  >{\raggedright\arraybackslash}p{(\columnwidth - 2\tabcolsep) * \real{0.5000}}@{}}
\toprule\noalign{}
\begin{minipage}[b]{\linewidth}\raggedright
Setting
\end{minipage} & \begin{minipage}[b]{\linewidth}\raggedright
Value
\end{minipage} \\
\midrule\noalign{}
\endhead
\bottomrule\noalign{}
\endlastfoot
Cities & A-B-C synthetic corridor \\
Inter-city distance & 100.0 km \\
Fixed stations & 2 (one between A-B and one between B-C) \\
Fixed plugs & 12 at AB, 12 at BC \\
Mobile stations & Up to 10 units \\
Chargers per MCS & 2 \\
Decision step & 5 minutes \\
Episode duration & 2-6 days during training \\
Mean service time & 30.0 minutes \\
Morning peak & 8.0:00 \\
Evening peak & 17.0:00 \\
Disruption mode & random \\
Supported disruption types & capacity\_drop, station\_outage,
service\_time\_inflation, demand\_surge \\
Mean MCS throughput & 4.0 vehicles/hour per MCS (\textasciitilde0.33
vehicles per 5-minute step) \\
Depot-to-station travel & 30 minutes \\
Full MCS recharge time & 60 minutes \\
\end{longtable}

    
    \begin{center}
    \adjustimage{max size={0.9\linewidth}{0.9\paperheight}}{exam_storyline_files/exam_storyline_4_1.png}
    \end{center}
    { \hspace*{\fill} \\}
    
    The schematic highlights the key spatial constraint in the project. The
depot is in the middle, while congestion may emerge on either side of
the corridor. A controller therefore faces a real allocation problem:
support sent to AB is not simultaneously available at BC, and vice
versa.

    \subsection{3. Simulation design}\label{simulation-design}

This section explains how the simulator is built up step by step. The
purpose is not to mimic every real-world detail, but to construct an
operationally meaningful resilience testbed in which the control problem
is interpretable.

\subsubsection{3.1 Base corridor}\label{base-corridor}

The starting point is the simplest spatial setting that still has local
congestion: three cities on a line and one fixed charging station on
each link. This gives two service points, AB and BC, that can become
stressed differently.

\subsubsection{3.2 Time-varying demand}\label{time-varying-demand}

The next layer is demand variation over the day. Morning and evening
peaks are added so that queue pressure is not constant. This matters
because a resilience intervention should help most when traffic is
concentrated, not only when the system is quiet.

\subsubsection{3.3 OD structure}\label{od-structure}

Demand is then split into origin-destination flows rather than a single
generic arrival stream. That matters because some trips pressure only
one station while long trips pressure both. The corridor therefore has
directional structure, which is important for spatial allocation.

\subsubsection{3.4 Queueing and service}\label{queueing-and-service}

Passing traffic is converted into charging demand. Some vehicles decide
to stop, they join the queue of the relevant station, and they occupy a
plug for a stochastic service time. This is the step that turns traffic
variation into a service-level problem with waiting times.

\subsubsection{3.5 Mobile charging
stations}\label{mobile-charging-stations}

Mobile charging stations are added only after the fixed charging system
exists. Their role is to add flexible capacity, not replace the fixed
system. Introducing travel and recharge delays makes the MCS decision
spatial and delayed rather than trivial.

\subsubsection{3.6 Disruptions}\label{disruptions}

Finally, disruption events are added. These reduce capacity, inflate
service time, or create temporary demand surges. This is what turns the
simulator into a resilience testbed rather than a normal-day capacity
exercise.

    \begin{longtable}[]{@{}
  >{\raggedright\arraybackslash}p{(\columnwidth - 4\tabcolsep) * \real{0.3333}}
  >{\raggedright\arraybackslash}p{(\columnwidth - 4\tabcolsep) * \real{0.3333}}
  >{\raggedright\arraybackslash}p{(\columnwidth - 4\tabcolsep) * \real{0.3333}}@{}}
\toprule\noalign{}
\begin{minipage}[b]{\linewidth}\raggedright
Layer added
\end{minipage} & \begin{minipage}[b]{\linewidth}\raggedright
Implementation
\end{minipage} & \begin{minipage}[b]{\linewidth}\raggedright
Why it matters
\end{minipage} \\
\midrule\noalign{}
\endhead
\bottomrule\noalign{}
\endlastfoot
Base corridor & Three cities in a line with one fixed station on each
link. & Creates the minimal spatial setting in which local congestion
can differ between AB and BC. \\
Time-varying demand & Morning peak, afternoon peak, and a smaller midday
bump. & Makes congestion depend on time of day rather than a flat
average flow. \\
OD structure & Trips are split across AB, BA, BC, CB, AC, and CA flows.
& Allows station pressure to come from directional travel patterns, not
only total volume. \\
Queueing and service & Vehicles may stop to charge, join queues, and
occupy plugs for stochastic service times. & Turns traffic into an
operational service problem with waiting-time consequences. \\
Mobile charging stations & Up to 10 MCS units add flexible but delayed
capacity. & Creates the control lever that can potentially improve
resilience. \\
Disruptions & Capacity drops, outages, service-time inflation, and OD
demand surges. & Creates stress conditions under which resilience can be
evaluated directly. \\
\end{longtable}

    
    \begin{center}
    \adjustimage{max size={0.9\linewidth}{0.9\paperheight}}{exam_storyline_files/exam_storyline_7_1.png}
    \end{center}
    { \hspace*{\fill} \\}
    
    The demand plot is the first important building block. The top panel
shows that expected charging demand is structured around a morning peak,
an afternoon/evening peak, and a smaller midday bump. The lower panel
shows the OD components that generate this pattern, which is why
congestion can become directional rather than purely aggregate.

    \begin{longtable}[]{@{}
  >{\raggedright\arraybackslash}p{(\columnwidth - 4\tabcolsep) * \real{0.3333}}
  >{\raggedright\arraybackslash}p{(\columnwidth - 4\tabcolsep) * \real{0.3333}}
  >{\raggedright\arraybackslash}p{(\columnwidth - 4\tabcolsep) * \real{0.3333}}@{}}
\toprule\noalign{}
\begin{minipage}[b]{\linewidth}\raggedright
Process
\end{minipage} & \begin{minipage}[b]{\linewidth}\raggedright
Distribution / mechanism
\end{minipage} & \begin{minipage}[b]{\linewidth}\raggedright
Assumption in the active benchmark
\end{minipage} \\
\midrule\noalign{}
\endhead
\bottomrule\noalign{}
\endlastfoot
Expected total traffic & Deterministic baseline + Gaussian-shaped
morning/evening peaks + midday bump & Baseline 26 cars/step; peak width
1.6 h. \\
Origin shares & Time-varying weighted shares across A, B, and C origins
& Directional bias shifts origin pressure toward A in the morning and
toward C in the evening. \\
Passing vehicles by OD and step & Poisson draw around expected OD flow &
Implements stochastic arrivals around the deterministic demand
profile. \\
Charging decision & Binomial draw from passing vehicles with fixed stop
probability & Charging-stop probability = 0.055. \\
Station assignment & Multinomial assignment using OD-specific station
weights & Long trips can pressure both stations; local trips pressure
one side only. \\
Service time & Normal draw, then clipped to configured min/max bounds &
Mean 30 min, std 8, clipped to 15-50 min. \\
Disruptions per day & Categorical draw from configured day-level count
weights & Weights: \{0: 0.1, 1: 0.45, 2: 0.3, 3: 0.15\}. \\
Disruption timing and severity & Uniform draws within configured ranges
for start time, duration, and severity & Keeps disruptions stochastic
but bounded and interpretable. \\
\end{longtable}

    
    \begin{longtable}[]{@{}
  >{\raggedright\arraybackslash}p{(\columnwidth - 8\tabcolsep) * \real{0.2000}}
  >{\raggedright\arraybackslash}p{(\columnwidth - 8\tabcolsep) * \real{0.2000}}
  >{\raggedright\arraybackslash}p{(\columnwidth - 8\tabcolsep) * \real{0.2000}}
  >{\raggedright\arraybackslash}p{(\columnwidth - 8\tabcolsep) * \real{0.2000}}
  >{\raggedright\arraybackslash}p{(\columnwidth - 8\tabcolsep) * \real{0.2000}}@{}}
\toprule\noalign{}
\begin{minipage}[b]{\linewidth}\raggedright
Type
\end{minipage} & \begin{minipage}[b]{\linewidth}\raggedright
Targets
\end{minipage} & \begin{minipage}[b]{\linewidth}\raggedright
Start window
\end{minipage} & \begin{minipage}[b]{\linewidth}\raggedright
Duration
\end{minipage} & \begin{minipage}[b]{\linewidth}\raggedright
Severity
\end{minipage} \\
\midrule\noalign{}
\endhead
\bottomrule\noalign{}
\endlastfoot
capacity drop & station\_ab, station\_bc & 6.0-20.0 h & {[}1.0, 3.5{]} &
Target plugs sampled from {[}1, 5{]} \\
station outage & station\_ab, station\_bc & 6.0-21.0 h & {[}1.0, 2.5{]}
& Target station forced to zero plugs \\
service time inflation & station\_ab, station\_bc & 6.0-20.0 h & {[}1.0,
3.5{]} & Service-time multiplier sampled from {[}1.5, 2.1{]} \\
demand surge & od\_ab, od\_ba, od\_cb, od\_bc, od\_ac, od\_ca & 6.0-20.0
h & {[}1.5, 4.5{]} & Demand multiplier sampled from {[}1.9, 2.9{]} \\
\end{longtable}

    
    Two design choices deserve emphasis.

First, the simulator mixes deterministic structure and stochastic
realization. The daily demand profile is interpretable, but arrivals and
service times are still random. Second, disruptions are sampled from
bounded ranges. That keeps the scenarios diverse enough to test
robustness, while still allowing the examiner to understand what a
disruption actually means in this benchmark.

The four subsections below switch from the full random benchmark to
\textbf{scripted one-day examples}. Each example contains exactly one
disruption type so its operational mechanism is easy to see. The goal is
not to produce benchmark scores here, but to show how each disruption
changes the relationship between expected demand, queue length, and
queue wait.

    \paragraph{3.6.1 Capacity drop}\label{capacity-drop}

A capacity-drop disruption reduces the number of usable fixed plugs at a
targeted station without changing the incoming demand. It therefore
represents a service-capability shock rather than a traffic shock.

    \begin{center}
    \adjustimage{max size={0.9\linewidth}{0.9\paperheight}}{exam_storyline_files/exam_storyline_12_0.png}
    \end{center}
    { \hspace*{\fill} \\}
    
    The expected-demand line remains on its normal path, while queue length
and waiting time rise sharply inside the shaded window. That is the
signature of a resilience problem driven by lost local service capacity
rather than extra arrivals.

    \paragraph{3.6.2 Station outage}\label{station-outage}

A station-outage disruption is the most severe local service failure in
the benchmark. The targeted station is forced completely offline for a
limited period, so all arriving charging demand must either wait or be
absorbed later when capacity returns.

    \begin{center}
    \adjustimage{max size={0.9\linewidth}{0.9\paperheight}}{exam_storyline_files/exam_storyline_15_0.png}
    \end{center}
    { \hspace*{\fill} \\}
    
    This case produces the steepest queue response because effective local
capacity falls all the way to zero. The demand path still does not move
much, which makes the disruption effect especially easy to separate from
ordinary morning congestion.

    \paragraph{3.6.3 Service-time inflation}\label{service-time-inflation}

A service-time inflation disruption does not remove plugs, but it makes
each charging session last longer. Operationally, this mimics slower
charging, faulty chargers operating below normal speed, or supporting
processes that delay turnover.

    \begin{center}
    \adjustimage{max size={0.9\linewidth}{0.9\paperheight}}{exam_storyline_files/exam_storyline_18_0.png}
    \end{center}
    { \hspace*{\fill} \\}
    
    The queue response is milder than a full outage but more persistent,
because congestion builds through slower turnover rather than an
immediate hard stop. This is a useful reminder that resilience problems
can come from degraded service rates as well as from missing plugs.

    \paragraph{3.6.4 Demand surge}\label{demand-surge}

A demand-surge disruption differs from the previous three because it
changes the arrival side of the system rather than the service side. In
the scripted example below, eastbound OD flows are multiplied during the
morning rush, which pushes more charging demand into the corridor before
any extra fixed capacity is available.

    \begin{center}
    \adjustimage{max size={0.9\linewidth}{0.9\paperheight}}{exam_storyline_files/exam_storyline_21_0.png}
    \end{center}
    { \hspace*{\fill} \\}
    
    This figure looks different from the previous three because the
expected-demand line itself increases during the disruption window. That
contrast is useful in the report: capacity-side disruptions create queue
spikes without extra demand, while demand surges raise congestion by
directly increasing the incoming load.

    \subsection{4. RL agent and methods}\label{rl-agent-and-methods}

The repository contains two RL paths:

\begin{itemize}
\tightlist
\item
  a \textbf{simple PyTorch DQN / DQL path} that writes full train/eval
  diagnostics,
\item
  an \textbf{SB3 DQN path} that is the intended stronger final policy
  path.
\end{itemize}

For the active benchmark, the action problem is deliberately posed as a
\textbf{small directional control problem} rather than a large
absolute-allocation table. The agent does not choose an arbitrary global
allocation in one jump. Instead, it chooses whether to hold, push
commitment toward AB or BC, or recall commitment from one side. This is
easier to learn, easier to interpret, and more faithful to the
operational story of gradual reallocation.

Because the current checkout does not contain a compatible final RL
checkpoint, this report distinguishes between the \textbf{RL design} and
the \textbf{reproducible final rollout evidence}. The RL design can
still be described precisely and evaluated through training diagnostics.

    \begin{longtable}[]{@{}
  >{\raggedright\arraybackslash}p{(\columnwidth - 2\tabcolsep) * \real{0.5000}}
  >{\raggedright\arraybackslash}p{(\columnwidth - 2\tabcolsep) * \real{0.5000}}@{}}
\toprule\noalign{}
\begin{minipage}[b]{\linewidth}\raggedright
Item
\end{minipage} & \begin{minipage}[b]{\linewidth}\raggedright
Value
\end{minipage} \\
\midrule\noalign{}
\endhead
\bottomrule\noalign{}
\endlastfoot
Observation size & 27 \\
Action size & 5 \\
Simple learner backend & torch\_dqn \\
Simple learner network & MLP 27 → 256 → 256 → 5 \\
Simple learner optimizer & Adam, learning rate 0.0005 \\
Simple learner TD loss & Huber / smooth L1 loss on one-step TD
targets \\
SB3 candidate backend & sb3\_dqn \\
SB3 policy class & MlpPolicy with net\_arch {[}256, 256{]} \\
Current checkpoint status &
exam\_submission/demo\_outputs/exam\_demo\_mobile\_mcs\_line\_abc/rl/best\_model.pt \\
\end{longtable}

    
    \begin{longtable}[]{@{}ll@{}}
\toprule\noalign{}
Simple learner hyperparameter & Value \\
\midrule\noalign{}
\endhead
\bottomrule\noalign{}
\endlastfoot
Episodes & 130 \\
Max steps / episode & 1728 \\
Discount factor γ & 0.98 \\
Learning rate & 0.0005 \\
Batch size & 128 \\
Replay capacity & 200000 \\
Learning starts & 10000 \\
Train frequency & 4 \\
Gradient steps & 4 \\
Target update interval & 200 \\
Hidden dimensions & {[}256, 256{]} \\
Exploration start ε & 1.0 \\
Exploration end ε & 0.05 \\
Exploration decay steps & 125000 \\
Periodic eval cadence & Every 15000 environment steps \\
Periodic eval episodes & 10 \\
\end{longtable}

    
    \begin{longtable}[]{@{}ll@{}}
\toprule\noalign{}
SB3 hyperparameter & Value \\
\midrule\noalign{}
\endhead
\bottomrule\noalign{}
\endlastfoot
Total timesteps & 250000.0 \\
Learning rate & 0.0005 \\
Discount factor γ & 0.98 \\
Replay buffer & 200000.0 \\
Learning starts & 10000.0 \\
Batch size & 128.0 \\
Train frequency & 4.0 \\
Gradient steps & 1.0 \\
Exploration fraction & 0.5 \\
Initial ε & 1.0 \\
Final ε & 0.05 \\
Target update interval & 1000.0 \\
Evaluation episodes & 12.0 \\
\end{longtable}

    
    \subsubsection{State space}\label{state-space}

The observation is 27-dimensional. The model is given local queue
pressure, recent flow, effective local capacity, mobile-station
logistics, disruption indicators, and time-of-day information. This
means the policy can react to where pressure is building, but it does
\textbf{not} have perfect future knowledge.

    \begin{longtable}[]{@{}
  >{\raggedright\arraybackslash}p{(\columnwidth - 4\tabcolsep) * \real{0.3333}}
  >{\raggedright\arraybackslash}p{(\columnwidth - 4\tabcolsep) * \real{0.3333}}
  >{\raggedright\arraybackslash}p{(\columnwidth - 4\tabcolsep) * \real{0.3333}}@{}}
\toprule\noalign{}
\begin{minipage}[b]{\linewidth}\raggedright
Feature group
\end{minipage} & \begin{minipage}[b]{\linewidth}\raggedright
Count
\end{minipage} & \begin{minipage}[b]{\linewidth}\raggedright
Why it is in the state
\end{minipage} \\
\midrule\noalign{}
\endhead
\bottomrule\noalign{}
\endlastfoot
Queue state & 4 & Current queue length and current mean wait at AB and
BC. \\
Recent flow & 4 & Latest arrivals and latest service starts at each
station. \\
Capacity state & 4 & Effective plugs and active mobile capacity at AB
and BC. \\
Spatial pressure & 4 & Committed MCS bias and local deficit estimates by
side of corridor. \\
Disruption state & 4 & Whether a disruption is active, what type it is,
and which side it affects. \\
MCS logistics & 5 & Units available in depot, charging, or in transit to
AB, BC, or middle. \\
Time-of-day & 2 & Sine and cosine encodings of time within the day. \\
\end{longtable}

    
    \begin{itemize}
\tightlist
\item
  No exact future arrivals or exact future queue realizations.
\end{itemize}

    
    \begin{itemize}
\tightlist
\item
  No oracle action labels or optimal allocation target.
\end{itemize}

    
    \begin{itemize}
\tightlist
\item
  No foreknowledge of future random disruptions before they begin.
\end{itemize}

    
    \begin{itemize}
\tightlist
\item
  No exact per-MCS battery state-of-charge representation.
\end{itemize}

    
    \begin{itemize}
\tightlist
\item
  No direct remaining-disruption-time feature in the current
  observation.
\end{itemize}

    
    \begin{itemize}
\tightlist
\item
  No explicit demand forecast beyond current local state and disruption
  indicators.
\end{itemize}

    
    \subsubsection{Action space}\label{action-space}

The active action space has five commands: hold, move one unit toward
AB, move one unit toward BC, recall one unit from AB, and recall one
unit from BC. This should be read as a deliberate modeling choice, not a
reduced convenience version. It keeps the controller focused on
directional decisions that respect depot availability and relocation
lag.

A subtle but important detail is that the environment itself decides
whether a \texttt{toward\_ab} command is satisfied from the middle depot
or by shifting one committed unit from BC. The policy chooses the
\textbf{directional intent}; the environment resolves the
\textbf{physical logistics}.

    \begin{longtable}[]{@{}
  >{\raggedright\arraybackslash}p{(\columnwidth - 2\tabcolsep) * \real{0.5000}}
  >{\raggedright\arraybackslash}p{(\columnwidth - 2\tabcolsep) * \real{0.5000}}@{}}
\toprule\noalign{}
\begin{minipage}[b]{\linewidth}\raggedright
Action
\end{minipage} & \begin{minipage}[b]{\linewidth}\raggedright
Interpretation
\end{minipage} \\
\midrule\noalign{}
\endhead
\bottomrule\noalign{}
\endlastfoot
hold & Keep the committed mobile-station allocation unchanged. \\
toward\_ab & Move one unit of commitment toward AB, either from depot or
from BC. \\
toward\_bc & Move one unit of commitment toward BC, either from depot or
from AB. \\
recall\_ab & Recall one committed unit away from AB toward the middle
depot. \\
recall\_bc & Recall one committed unit away from BC toward the middle
depot. \\
\end{longtable}

    
    \subsubsection{Reward function and loss
function}\label{reward-function-and-loss-function}

The reward combines service-quality incentives and operating-cost
penalties.

The positive terms reward the agent for serving vehicles, using capacity
productively, and placing mobile support where local deficits are
estimated to be highest. The negative terms penalize queueing, waiting,
unmet demand, unnecessary MCS activation, repeated adjustment, and
carrying idle capacity.

This distinction matters because \textbf{operational performance} and
\textbf{reward} are related but not identical. A policy may reduce
queueing substantially while still earning a poor total reward if it
uses too much costly mobile capacity.

The simple learner is trained with a \textbf{Huber TD loss}
(\texttt{smooth\_l1\_loss}). Intuitively, the model predicts action
values, compares them with one-step bootstrap targets, and updates the
network to reduce that temporal-difference error. Huber loss is used
because it is less brittle than pure squared error when RL targets are
noisy.

    \begin{longtable}[]{@{}
  >{\raggedright\arraybackslash}p{(\columnwidth - 4\tabcolsep) * \real{0.3333}}
  >{\raggedright\arraybackslash}p{(\columnwidth - 4\tabcolsep) * \real{0.3333}}
  >{\raggedright\arraybackslash}p{(\columnwidth - 4\tabcolsep) * \real{0.3333}}@{}}
\toprule\noalign{}
\begin{minipage}[b]{\linewidth}\raggedright
Positive term
\end{minipage} & \begin{minipage}[b]{\linewidth}\raggedright
Weight
\end{minipage} & \begin{minipage}[b]{\linewidth}\raggedright
Operational role
\end{minipage} \\
\midrule\noalign{}
\endhead
\bottomrule\noalign{}
\endlastfoot
served\_reward\_weight & 1.0 & Rewards charging sessions that start
service. \\
utilization\_bonus & 1.0 & Rewards productive use of available charging
capacity. \\
spatial\_deficit\_alignment\_bonus & 24.0 & Rewards mobile capacity that
covers estimated local deficits. \\
spatial\_deficit\_direction\_bonus & 12.0 & Rewards sending MCS in the
correct corridor direction. \\
\end{longtable}

    
    \begin{longtable}[]{@{}
  >{\raggedright\arraybackslash}p{(\columnwidth - 4\tabcolsep) * \real{0.3333}}
  >{\raggedright\arraybackslash}p{(\columnwidth - 4\tabcolsep) * \real{0.3333}}
  >{\raggedright\arraybackslash}p{(\columnwidth - 4\tabcolsep) * \real{0.3333}}@{}}
\toprule\noalign{}
\begin{minipage}[b]{\linewidth}\raggedright
Negative term
\end{minipage} & \begin{minipage}[b]{\linewidth}\raggedright
Weight
\end{minipage} & \begin{minipage}[b]{\linewidth}\raggedright
Operational role
\end{minipage} \\
\midrule\noalign{}
\endhead
\bottomrule\noalign{}
\endlastfoot
unmet\_penalty & 2.5 & Penalizes vehicles left in queue after the
step. \\
queue\_length\_penalty & 2.5 & Adds extra pressure against persistent
queue accumulation. \\
queue\_wait\_penalty & 0.02 & Penalizes queue-wait burden, not just
queue count. \\
local\_queue\_peak\_penalty & 1.0 & Penalizes the worst station queue in
the step. \\
local\_wait\_peak\_penalty & 0.1 & Penalizes the worst station waiting
time in the step. \\
active\_mobile\_station\_cost & 18.0 & Charges for keeping MCS active in
the field. \\
activation\_cost & 2.0 & Charges for newly activating MCS. \\
adjustment\_cost & 0.5 & Charges for reallocating MCS commitment. \\
idle\_capacity\_penalty & 0.1 & Penalizes carrying unused capacity. \\
\end{longtable}

    
    \subsection{5. Reward calibration}\label{reward-calibration}

Reward calibration is a core modeling choice in this project. Mobile
charging can reduce queues, but it also represents operating cost. If
the MCS cost weight is too low, the controller may over-deploy mobile
support. If the queue penalty is too low, the controller may tolerate
poor service quality instead of protecting the queue.

The historical sweep available in the repository is sparse rather than a
full grid. To keep the section clean and defensible, the two plots below
use only the \textbf{one-factor slices} where one reward weight changes
while the other is held fixed.

    \begin{center}
    \adjustimage{max size={0.9\linewidth}{0.9\paperheight}}{exam_storyline_files/exam_storyline_32_0.png}
    \end{center}
    { \hspace*{\fill} \\}
    
    \begin{longtable}[]{@{}lllll@{}}
\toprule\noalign{}
Run & MCS cost & Queue cost (fixed) & Mean queue & Final reward \\
\midrule\noalign{}
\endhead
\bottomrule\noalign{}
\endlastfoot
range\_c18\_q2p5 & 18.0 & 2.5 & 0.62 & -102.93 \\
range\_c24\_q2p5 & 24.0 & 2.5 & 0.12 & -146.77 \\
\end{longtable}

    
    \begin{center}
    \adjustimage{max size={0.9\linewidth}{0.9\paperheight}}{exam_storyline_files/exam_storyline_32_2.png}
    \end{center}
    { \hspace*{\fill} \\}
    
    \begin{longtable}[]{@{}lllll@{}}
\toprule\noalign{}
Run & Queue cost & MCS cost (fixed) & Mean queue & Final reward \\
\midrule\noalign{}
\endhead
\bottomrule\noalign{}
\endlastfoot
range\_c30\_q3 & 3.0 & 30.0 & 2.25 & -232.35 \\
range\_c30\_q4 & 4.0 & 30.0 & 2.0 & -199.99 \\
\end{longtable}

    
    These plots should be read as simple directional sensitivity checks
rather than as a full optimization study.

\begin{itemize}
\tightlist
\item
  The \textbf{MCS-cost plot} isolates what happens when mobile support
  becomes more expensive while the queue penalty stays fixed.
\item
  The \textbf{queue-cost plot} isolates what happens when poor queue
  performance becomes more expensive while the MCS cost stays fixed.
\item
  The small tables below each plot show the corresponding reward
  settings and the resulting final reward, which helps connect the queue
  outcome back to the reward function that produced it.
\end{itemize}

The broader modeling point is that there is no single universally
correct reward. The preferred calibration depends on the operational
question: if mobile charging is expensive, conservative deployment may
be rational; if queueing is very costly, the controller should be
willing to spend more to protect service quality.

    \subsection{6. Training setup}\label{training-setup}

The training pipeline separates training episodes, periodic evaluation,
held-out rollouts, and named stress scenarios.

\begin{itemize}
\tightlist
\item
  \textbf{Training} uses randomized episodes from the benchmark
  distribution.
\item
  \textbf{Periodic evaluation} checks the current policy on held-out
  seeds without exploration and without learning updates.
\item
  The \textbf{held-out rollout} is the final time-series comparison used
  in the main report.
\item
  \textbf{Stress scenarios} are scripted one-day disruptions used as
  robustness checks and are kept in the appendix.
\end{itemize}

Even without hyperparameter tuning or early stopping, periodic
evaluation is still useful. It shows whether the policy is improving on
held-out scenarios rather than only fitting the randomness of the latest
training episodes.

    \begin{longtable}[]{@{}
  >{\raggedright\arraybackslash}p{(\columnwidth - 2\tabcolsep) * \real{0.5000}}
  >{\raggedright\arraybackslash}p{(\columnwidth - 2\tabcolsep) * \real{0.5000}}@{}}
\toprule\noalign{}
\begin{minipage}[b]{\linewidth}\raggedright
Component
\end{minipage} & \begin{minipage}[b]{\linewidth}\raggedright
How it is used in this report
\end{minipage} \\
\midrule\noalign{}
\endhead
\bottomrule\noalign{}
\endlastfoot
Training episodes & Randomized 2-6 day episodes sampled from training
seeds 0-9999. \\
Periodic evaluation & 10 held-out episodes every 15000 environment steps
without exploration or learning updates. \\
Validation split & Configured held-out validation seeds starting at
10000; useful for policy checks even without hyperparameter tuning. \\
Held-out rollout & A fixed unseen \texttt{test\_id} scenario used for
the main time-series comparison in this report. \\
Test stress suite & Named one-day scripted disruptions used for
robustness checks; kept in the appendix to avoid clutter. \\
\end{longtable}

    
    \subsection{7. Training the model}\label{training-the-model}

The final intended model path is SB3 DQN. However, the current
repository snapshot does not contain a compatible checkpoint for the
active benchmark, so the notebook cannot load and roll out that policy
directly.

What \textbf{is} available is a copied current training-history
artifact. That artifact contains train reward, train TD loss, evaluation
reward, and evaluation TD loss. These diagnostics are therefore used as
evidence about whether the RL learner improved during training.

The optional demo cell below is intentionally separate from the reported
results. Its only purpose is to show how an examiner can launch a short
local train/eval run using the included configs and reference code.

    \begin{longtable}[]{@{}
  >{\raggedright\arraybackslash}p{(\columnwidth - 2\tabcolsep) * \real{0.5000}}
  >{\raggedright\arraybackslash}p{(\columnwidth - 2\tabcolsep) * \real{0.5000}}@{}}
\toprule\noalign{}
\begin{minipage}[b]{\linewidth}\raggedright
Artifact check
\end{minipage} & \begin{minipage}[b]{\linewidth}\raggedright
Status
\end{minipage} \\
\midrule\noalign{}
\endhead
\bottomrule\noalign{}
\endlastfoot
Training-history source & Copied current training history in
exam\_submission/data/generated/training/history\_current.json \\
Train reward available & Yes \\
Train TD loss available & Yes \\
Eval reward available & Yes \\
Eval TD loss available & Yes \\
Current compatible checkpoint available & Yes \\
\end{longtable}

    
    \begin{longtable}[]{@{}
  >{\raggedright\arraybackslash}p{(\columnwidth - 2\tabcolsep) * \real{0.5000}}
  >{\raggedright\arraybackslash}p{(\columnwidth - 2\tabcolsep) * \real{0.5000}}@{}}
\toprule\noalign{}
\begin{minipage}[b]{\linewidth}\raggedright
Checkpoint
\end{minipage} & \begin{minipage}[b]{\linewidth}\raggedright
Exists
\end{minipage} \\
\midrule\noalign{}
\endhead
\bottomrule\noalign{}
\endlastfoot
exam\_submission/demo\_outputs/exam\_demo\_mobile\_mcs\_line\_abc/rl/best\_model.pt
& True \\
\end{longtable}

    
    \begin{center}
    \adjustimage{max size={0.9\linewidth}{0.9\paperheight}}{exam_storyline_files/exam_storyline_37_2.png}
    \end{center}
    { \hspace*{\fill} \\}
    
    The learning curves should be interpreted carefully.

\begin{itemize}
\tightlist
\item
  Train reward and TD loss are noisy, as expected in RL.
\item
  Evaluation reward is more informative than raw training reward because
  it is measured without exploration.
\item
  Evaluation TD loss is not a supervised validation loss, but it is
  still useful as a consistency check on the learned value function.
\end{itemize}

The copied history suggests that the learner improved relative to its
very poor early episodes and reached a more stable regime. However,
because the compatible final checkpoint is missing, the report cannot
make a final reproducible claim about RL rollout performance on the
active benchmark.

    \subsubsection{Optional demo run}\label{optional-demo-run}

The demo run is disabled by default so the notebook stays fast and
deterministic. To run a short local train/eval demonstration, set
\texttt{RUN\_DEMO\_TRAINING\ =\ False} in this cell.

\begin{Shaded}
\begin{Highlighting}[]
\ExtensionTok{/Library/Frameworks/Python.framework/Versions/3.11/bin/python3.11} \AttributeTok{{-}m}\NormalTok{ evch.train.train\_rl }\AttributeTok{{-}{-}config}\NormalTok{ /Users/nicolaigarderhansen/Desktop/DTU/Kandidat/2. Sem/42578/advanced{-}business{-}analytics/exam\_submission/reference\_code/configs/env/mobile\_mcs\_line\_abc.yaml }\AttributeTok{{-}{-}config}\NormalTok{ /Users/nicolaigarderhansen/Desktop/DTU/Kandidat/2. Sem/42578/advanced{-}business{-}analytics/exam\_submission/reference\_code/configs/demand/base.yaml }\AttributeTok{{-}{-}config}\NormalTok{ /Users/nicolaigarderhansen/Desktop/DTU/Kandidat/2. Sem/42578/advanced{-}business{-}analytics/exam\_submission/reference\_code/configs/rl/dqn\_mobile\_simple.yaml }\AttributeTok{{-}{-}config}\NormalTok{ /Users/nicolaigarderhansen/Desktop/DTU/Kandidat/2. Sem/42578/advanced{-}business{-}analytics/exam\_submission/reference\_code/configs/debug/logging\_disabled.yaml }\AttributeTok{{-}{-}config}\NormalTok{ /Users/nicolaigarderhansen/Desktop/DTU/Kandidat/2. Sem/42578/advanced{-}business{-}analytics/exam\_submission/reference\_code/configs/experiment/mobile\_mcs\_line\_abc\_train\_only.yaml }\AttributeTok{{-}{-}config}\NormalTok{ /Users/nicolaigarderhansen/Desktop/DTU/Kandidat/2. Sem/42578/advanced{-}business{-}analytics/exam\_submission/reference\_code/configs/demo/mobile\_mcs\_line\_abc\_demo\_short.yaml}
\end{Highlighting}
\end{Shaded}

    
    \subsection{8. Simulation comparison
rollout}\label{simulation-comparison-rollout}

The intended main comparison was fixed-only capacity versus the final RL
controller. Because no current-compatible RL rollout artifact is present
in the checkout, that exact comparison cannot be reproduced honestly in
this notebook.

To keep the report reproducible, the main operational comparison below
uses:

\begin{itemize}
\tightlist
\item
  the \textbf{fixed / no-agent baseline}, and
\item
  the \textbf{adaptive threshold controller}, which is the strongest
  current adaptive rollout artifact available.
\end{itemize}

This still answers an important resilience question: does adaptive
mobile support materially reduce queueing pressure relative to doing
nothing adaptive at all? The appendix retains the reactive rule as
secondary context, but it is removed from the main section to keep the
story clean.

    \begin{center}
    \adjustimage{max size={0.9\linewidth}{0.9\paperheight}}{exam_storyline_files/exam_storyline_41_0.png}
    \end{center}
    { \hspace*{\fill} \\}
    
    The held-out rollout shows a clear operational improvement. Fixed
capacity alone allows queue length and waiting time to build up
substantially during disruption windows. The adaptive threshold
controller responds by activating mobile support, which reduces both
total queueing and average waiting time.

The plot also illustrates the cost side of the problem: queue
improvement is achieved by using MCS, not by creating capacity from
nowhere. That is why the reward section and the reward-calibration
section matter for interpretation.

    \subsection{9. Summary metrics table}\label{summary-metrics-table}

The table below summarizes the held-out rollout used in the main
comparison. Breach-rate columns are shown for completeness, but they are
not informative in the active exported artifacts because the queue-wait
target is zero throughout the rollout files.

    \begin{longtable}[]{@{}
  >{\raggedright\arraybackslash}p{(\columnwidth - 24\tabcolsep) * \real{0.0769}}
  >{\raggedright\arraybackslash}p{(\columnwidth - 24\tabcolsep) * \real{0.0769}}
  >{\raggedright\arraybackslash}p{(\columnwidth - 24\tabcolsep) * \real{0.0769}}
  >{\raggedright\arraybackslash}p{(\columnwidth - 24\tabcolsep) * \real{0.0769}}
  >{\raggedright\arraybackslash}p{(\columnwidth - 24\tabcolsep) * \real{0.0769}}
  >{\raggedright\arraybackslash}p{(\columnwidth - 24\tabcolsep) * \real{0.0769}}
  >{\raggedright\arraybackslash}p{(\columnwidth - 24\tabcolsep) * \real{0.0769}}
  >{\raggedright\arraybackslash}p{(\columnwidth - 24\tabcolsep) * \real{0.0769}}
  >{\raggedright\arraybackslash}p{(\columnwidth - 24\tabcolsep) * \real{0.0769}}
  >{\raggedright\arraybackslash}p{(\columnwidth - 24\tabcolsep) * \real{0.0769}}
  >{\raggedright\arraybackslash}p{(\columnwidth - 24\tabcolsep) * \real{0.0769}}
  >{\raggedright\arraybackslash}p{(\columnwidth - 24\tabcolsep) * \real{0.0769}}
  >{\raggedright\arraybackslash}p{(\columnwidth - 24\tabcolsep) * \real{0.0769}}@{}}
\toprule\noalign{}
\begin{minipage}[b]{\linewidth}\raggedright
Policy
\end{minipage} & \begin{minipage}[b]{\linewidth}\raggedright
Mean queue
\end{minipage} & \begin{minipage}[b]{\linewidth}\raggedright
Peak queue
\end{minipage} & \begin{minipage}[b]{\linewidth}\raggedright
Mean wait (min)
\end{minipage} & \begin{minipage}[b]{\linewidth}\raggedright
Peak wait (min)
\end{minipage} & \begin{minipage}[b]{\linewidth}\raggedright
Breach rate
\end{minipage} & \begin{minipage}[b]{\linewidth}\raggedright
Started sessions
\end{minipage} & \begin{minipage}[b]{\linewidth}\raggedright
Mean active MCS
\end{minipage} & \begin{minipage}[b]{\linewidth}\raggedright
MCS hours
\end{minipage} & \begin{minipage}[b]{\linewidth}\raggedright
Reward
\end{minipage} & \begin{minipage}[b]{\linewidth}\raggedright
Queue improvement vs fixed (\%)
\end{minipage} & \begin{minipage}[b]{\linewidth}\raggedright
Wait improvement vs fixed (\%)
\end{minipage} & \begin{minipage}[b]{\linewidth}\raggedright
Reward change vs fixed
\end{minipage} \\
\midrule\noalign{}
\endhead
\bottomrule\noalign{}
\endlastfoot
Fixed / no-agent & 16.6 & 97.0 & 12.37 & 78.14 & 0.0 & 3851.0 & 0.0 &
0.0 & -15969.37 & 0.0 & 0.0 & 0.0 \\
Adaptive threshold & 2.65 & 24.0 & 1.28 & 22.69 & 0.0 & 3846.0 & 0.95 &
113.5 & -4202.84 & 84.06 & 89.62 & 11766.53 \\
\end{longtable}

    
    The key result is that adaptive mobile support meaningfully reduces
queueing pain relative to fixed-only operation. The threshold controller
improves service outcomes while using much less mobile support than the
more aggressive reactive rule shown later in the appendix. That makes it
the cleanest main comparator among the currently reproducible adaptive
artifacts.

    \subsection{10. Spatial awareness}\label{spatial-awareness}

A spatially aware controller should allocate support where local
pressure is highest. In this corridor, that means MCS at AB should
increase when AB has higher local demand or queue pressure, while MCS at
BC should increase when BC is more stressed.

The threshold controller is the clearest current artifact for checking
this. The plots below show one-hour rolling averages of expected station
demand, local queue length, and MCS allocation for AB and BC separately.

    \begin{center}
    \adjustimage{max size={0.9\linewidth}{0.9\paperheight}}{exam_storyline_files/exam_storyline_47_0.png}
    \end{center}
    { \hspace*{\fill} \\}
    
    The spatial response is meaningful but not perfect.

\begin{itemize}
\tightlist
\item
  When local pressure rises, MCS allocation generally rises on the same
  side of the corridor.
\item
  The response is not instantaneous, which is expected because MCS units
  must travel and recharge.
\item
  The relationship is clearer for the threshold controller than for the
  reactive rule, because the reactive policy tends to keep many units
  active system-wide rather than expressing a clean directional
  response.
\end{itemize}

So the appropriate claim is not ``perfect spatial awareness''. A more
defensible conclusion is that the current simulator can reveal spatial
behavior clearly, and the adaptive threshold artifact shows a visible
directional response to local service pressure.

    \subsection{11. Limitations}\label{limitations}

Several limitations are important and should be stated directly.

\begin{enumerate}
\def\labelenumi{\arabic{enumi}.}
\tightlist
\item
  \textbf{The simulator is synthetic.} It is designed to be
  operationally meaningful, but it is not a calibrated production model
  of a real EV corridor.
\item
  \textbf{The current checkout lacks a compatible RL checkpoint.} The RL
  formulation and training diagnostics are available, but a final
  reproducible RL-vs-baseline rollout claim cannot be made from this
  repository snapshot alone.
\item
  \textbf{Reward sensitivity is substantial.} Different cost weights can
  change the preferred behavior significantly, especially the
  willingness to keep MCS active.
\item
  \textbf{Spatial behavior is partial rather than perfect.} Travel
  delay, recharge delay, directional symmetry, and cost all weaken
  one-to-one alignment between local pressure and instantaneous
  allocation.
\item
  \textbf{Logistics are simplified.} The simulator captures relocation
  and recharge delays, but not the full operational complexity of
  staffing, road-network uncertainty, or battery degradation.
\item
  \textbf{Stress coverage is limited.} The appendix shows several
  scripted stress scenarios, but this is still a bounded test set rather
  than a full uncertainty analysis over all possible disruptions.
\end{enumerate}

These limitations do not invalidate the project, but they do constrain
what can be claimed. The report therefore focuses on what is actually
supported by the copied artifacts.

    \subsection{12. Conclusion}\label{conclusion}

The project builds a coherent simulator for resilience in EV charging
under disruption on a simple but informative A-B-C corridor. The
environment combines time-varying demand, OD structure, queueing,
service dynamics, mobile charging logistics, and several disruption
types in a way that is interpretable from top to bottom.

Within the \textbf{current reproducible artifacts}, adaptive mobile
support clearly improves service outcomes relative to fixed-only
operation. The threshold controller reduces mean and peak queueing while
using a moderate amount of mobile support. The RL formulation itself is
well specified, and the copied train/eval curves suggest learning
progress, but the absence of a compatible current RL checkpoint prevents
a final reproducible RL rollout claim.

The most defensible final takeaway is therefore twofold: the simulator
is strong enough to study resilience and spatial allocation, and
adaptive mobile charging has clear value in that simulator. The next
concrete step would be to recover or retrain a current-compatible RL
checkpoint and benchmark it directly against the threshold baseline
documented here.

    \subsection{13. Appendix}\label{appendix}

The appendix keeps secondary material out of the main narrative while
still documenting the broader evidence base used in the project.

    \subsubsection{A. Full 27-feature observation
table}\label{a.-full-27-feature-observation-table}

    \begin{longtable}[]{@{}
  >{\raggedright\arraybackslash}p{(\columnwidth - 6\tabcolsep) * \real{0.2500}}
  >{\raggedright\arraybackslash}p{(\columnwidth - 6\tabcolsep) * \real{0.2500}}
  >{\raggedright\arraybackslash}p{(\columnwidth - 6\tabcolsep) * \real{0.2500}}
  >{\raggedright\arraybackslash}p{(\columnwidth - 6\tabcolsep) * \real{0.2500}}@{}}
\toprule\noalign{}
\begin{minipage}[b]{\linewidth}\raggedright
Group
\end{minipage} & \begin{minipage}[b]{\linewidth}\raggedright
Idx
\end{minipage} & \begin{minipage}[b]{\linewidth}\raggedright
Feature
\end{minipage} & \begin{minipage}[b]{\linewidth}\raggedright
Meaning
\end{minipage} \\
\midrule\noalign{}
\endhead
\bottomrule\noalign{}
\endlastfoot
Queue state & 1 & queue\_length\_ab & Current queue length at station
AB. \\
Queue state & 2 & queue\_length\_bc & Current queue length at station
BC. \\
Queue state & 3 & mean\_queue\_wait\_ab & Mean current waiting time in
the AB queue. \\
Queue state & 4 & mean\_queue\_wait\_bc & Mean current waiting time in
the BC queue. \\
Recent flow & 5 & last\_arrivals\_ab & Vehicles that arrived to AB in
the last step. \\
Recent flow & 6 & last\_arrivals\_bc & Vehicles that arrived to BC in
the last step. \\
Recent flow & 7 & last\_starts\_ab & Charging sessions started at AB in
the last step. \\
Recent flow & 8 & last\_starts\_bc & Charging sessions started at BC in
the last step. \\
Capacity & 9 & effective\_plugs\_ab & Effective plugs at AB after
disruption and MCS effects. \\
Capacity & 10 & effective\_plugs\_bc & Effective plugs at BC after
disruption and MCS effects. \\
Capacity & 11 & active\_mobile\_capacity\_ab & Currently active
mobile-station capacity at AB. \\
Capacity & 12 & active\_mobile\_capacity\_bc & Currently active
mobile-station capacity at BC. \\
Spatial & 13 & committed\_mcs\_bias\_ab\_minus\_bc & Committed MCS bias
toward AB versus BC. \\
Spatial & 14 & local\_deficit\_ab & Local deficit estimate at AB. \\
Spatial & 15 & local\_deficit\_bc & Local deficit estimate at BC. \\
Spatial & 16 & local\_deficit\_bias\_ab\_minus\_bc & Local deficit bias
toward AB versus BC. \\
Disruption & 17 & disruption\_active & Whether a disruption is currently
active. \\
Disruption & 18 & disruption\_type\_code & Encoded disruption type. \\
Disruption & 19 & target\_affects\_ab & Whether the disruption affects
AB. \\
Disruption & 20 & target\_affects\_bc & Whether the disruption affects
BC. \\
MCS logistics & 21 & middle\_available & Mobile stations available at
the middle depot. \\
MCS logistics & 22 & middle\_charging & Mobile stations currently
recharging at the middle depot. \\
MCS logistics & 23 & transit\_to\_ab & Mobile stations moving toward
AB. \\
MCS logistics & 24 & transit\_to\_bc & Mobile stations moving toward
BC. \\
MCS logistics & 25 & transit\_to\_middle & Mobile stations returning to
the middle depot. \\
Time & 26 & sin\_time\_of\_day & Sine encoding of time of day. \\
Time & 27 & cos\_time\_of\_day & Cosine encoding of time of day. \\
\end{longtable}

    

    % Add a bibliography block to the postdoc
    
    
    
\end{document}
